\cotefile{156-6}
\scanpage{74}\mspage{73}
J'ai envie de \unsure{renforcer} At $11'$, dans la direction suivante :

\uline{Proposition} (?) Soient $F,G\in\unsure{\mathcal F|\mathscr F}$, $F\mathrel{>\mkern-14mu<}G$, et soit $\deleted{\unsure{H}}K=F\cap G$. Soient $F'\ll F$, $G'\ll G$. Pour que $F'\mathrel{>\mkern-14mu<}G'$, il f. et s. que $F'|\deleted{\unsure{H}}K\deleted{\unsure{=}}\mathrel{>\mkern-14mu<}G'|\deleted{\unsure{H}}K$, et alors $F'\vee G'\ll F\vee G$.

Supposons que $F'$ et $G'$ soient compatibles, i.e. que $H'=F'\vee G'$ existe. \inserted{\deleted{Soit $K'=F'\wedge G'$}} Je dis que $H'\ll F\vee G\overset{\unsure{\text{défn}|\text{dfn}}}{=}H$ : cela résulte aussitôt du critère At 4. \deleted{Soit $K=F'\wedge G'$. On} \inserted{D'autre part, $F'|K$ étant une sous fig. de $F'$ est sous figure de $H'$, et de $=$ pour $G'|K$, donc} \deleted{$K'\deleted{\unsure{\ll}}\leq F'\deleted{\unsure{\leq}}\ll F$ \unsure{\&} et $K'\leq G'\ll G$, donc $K'\ll F,G$, donc ($\unsure{-\,-|\ldots}$) $K'\ll K=F\cap G=F.G$. Donc on a un diagramme cubique} \inserted{Ils sont \unsure{compatibles}. Inversement, supposons $F'|K\unsure{'}$ \deleted{\unsure{compati}} et $G'|K$ \unsure{compatibles}}

\diagram{\begin{tikzcd}[ampersand replacement=\&,row sep=0.45em,column sep=0.75em,execute at end picture={
\draw ([xshift=-3pt]pFa.south west) -- ([xshift=5pt,yshift=4pt]pF.north east);
\draw ([xshift=-6pt,yshift=-2pt]pKa.south west) -- ([xshift=4pt,yshift=4pt]pG.north west);
\draw ([yshift=-1pt]pKa.south east) -- ([xshift=6pt]pG.east);
}]
\& \& |[alias=pH]| \deleted{H} \& \\
|[alias=pF]| \deleted{F} \arrow[rru,hook] \& \& \& |[alias=pG]| \deleted{G} \arrow[ul,hook'] \\
\& |[alias=pK]| \deleted{K} \arrow[ul,leftrightarrow] \arrow[urr,hook] \& \& \\
\& \& \deleted{{\unsure{H}}'} \arrow[uuu] \& \\
|[alias=pFa]| \deleted{F'} \arrow[uuu] \arrow[rru,hook] \& \& \& \deleted{G'} \arrow[uuu] \arrow[ul,hook'] \\
\& |[alias=pKa]| \deleted{K'} \arrow[uuu] \arrow[ul,hook'] \arrow[urr,hook] \& \&
\end{tikzcd}}

\deleted{où les 8 flèches horizontales sont des $\leq$, \inserted{(les deux \unsure{faces|f.ces} horizontales sont cartésiennes)} et les quatre flèches verticales des $\ll$ sans plus.}

\deleted{On a}
\[
\begin{aligned}
&\deleted{K'=\deleted{F'\unsure{\mathord{.}|\cap}G'=F'}\operatorname{Inf}^{\ll}(F',G')\unsure{-}}\\
&\deleted{=\operatorname{Inf}^{\ll}(\deleted{\unsure{G|(}}F',G',K)=\operatorname{Inf}^{\ll}\bigl(\underbrace{\operatorname{Inf}(F',K)}_{F'|K},\underbrace{\operatorname{Inf}(G',K)}_{G'|K}\bigr)}
\end{aligned}
\]
\deleted{et $F'|K$ étant une sous-figure de $F'\leq H'$, en est une de $H'$, et de $=$ pour $G'$}
\scanpage{75}\mspage{74}
On aura \unsure{donc} un diagramme

\diagram{\begin{tikzcd}[ampersand replacement=\&,row sep=0.4em,column sep=0.85em]
\deleted{\unsure{F|E}} \& \& \& \& \\
\& \& H \& \& \\
F \arrow[rru,hook] \& \& \& G \arrow[ul,hook'] \& \\
\& K \arrow[ul,hook'] \arrow[urr,hook] \& \& \& \unsure{\rightharpoonup|\searrow} \\
\& \& \deleted{\unsure{H}}H' \arrow[uuu] \& \& \\
F' \arrow[uuu] \arrow[rru,hook] \& \& \& G' \arrow[uuu] \arrow[ul,hook'] \& \\
\& K' \arrow[uuu,two heads] \arrow[ul,hook'] \arrow[urr] \& \& \&
\end{tikzcd}}
\margin{\deleted{où $K'=F'\cap G'$}}

où $K'=F'\cap G'$. (NB. $K'\ll K$ car $K'\ll F$ et $K'\ll G$ \unsure{or|et} $K=\operatorname{Inf}^{\ll}(F,G)$) \deleted{On aura} Mieux,
\[
\begin{aligned}
&\deleted{K'=\operatorname{Inf}^{\ll}(F',G')=\operatorname{Inf}^{\ll}(F',G',K)=\operatorname{Inf}^{\ll}\bigl(\operatorname{Inf}^{\ll}(F',K),}\\
&\hspace{13em}\deleted{\operatorname{Inf}^{\ll}(G',\unsure{K|K^{\scriptscriptstyle ?}})\bigr)}
\end{aligned}
\]
\deleted{i.e.}
\[
\deleted{K'=\operatorname{Inf}\bigl(\underbrace{F'|K,\ G'|K}_{\substack{\text{deux \deleted{sous-}}\\\text{figures de }H'}}\bigr)
=(\underbrace{F'|K}_{F'_K})\cap(\underbrace{G'|K}_{G'_K}).}
\]
\deleted{et je dis que} \deleted{\unsure{et|A}}
\[
\begin{aligned}
&\deleted{\underbrace{H'|\deleted{\unsure{K}}F}_{H'_F}=F'\vee G'_K,\qquad H'|G=G'\vee F'_K}\\
&\qquad\deleted{\deleted{\unsure{=}}=F'\vee H'_K}\qquad\qquad\deleted{=G'\vee H'_K}
\end{aligned}
\]

\unsure{ou aussi|on aura}

\diagram{\begin{tikzcd}[ampersand replacement=\&,row sep=0.55em,column sep=0.6em]
\& \& \& H \& \& \\
\& F \arrow[rru,hook] \& \& \& G \arrow[ul,hook'] \& \\
\& \& K \arrow[ul,hook'] \arrow[urr,hook,no head] \& \& \& \\
\& \& \& H' \arrow[uuu] \& \& \\
\& H'_F \arrow[uuu] \arrow[rru,hook] \& \& \& \deleted{\unsure{H}}H'_G \arrow[uuu] \arrow[ul,hook'] \& \\
F' \arrow[ur,hook] \& \& H'_K \arrow[uuu] \arrow[ul,hook'] \arrow[urr,hook] \& \& \& G' \arrow[ul,hook'] \\
\& F'_K \arrow[ul,hook'] \arrow[ur,hook] \& \& \& G'_K \arrow[ull,hook'] \arrow[ur,hook] \& \\
\& \& K' \arrow[ul,hook'] \arrow[urr,hook] \& \& \&
\end{tikzcd}}
\margin{Diagramme où les cinq carrés horizontaux sont \deleted{\unsure{cartési} des $\leq$} cartésiens et cocartésiens (avec flèches $\xrightarrow{\leq}$) et les \deleted{\unsure{flè}} quatre flèches verticales sont des $\xrightarrow{\ll}$ sans plus.}
\scanpage{76}\mspage{75}
Inversement, supposons que $F'\ll F$ et $G'\ll G$ soient tels que $F'_K\mathrel{>\mkern-14mu<}G'_K$ \deleted{\unsure{soient|soit}, désignons par $H'_{\unsure{\mathord{.}|\mathord{'}}}$} \inserted{posons}
\[
\begin{aligned}
H'_K&=F'_K\vee G'_K &&\text{on aura }H'_K\ll K\\[-.3ex]
&&&\hspace{3em}\text{par At 4}\\
K'&=\underbrace{F'_K\cap G'_K}_{\unsure{\mathord{?}|\mathord{\&}}}\deleted{=F'\cap G'}
\end{aligned}
\]

Considérons $F'$ et $H'_K$ \deleted{\unsure{on voit} aussitôt que} \inserted{\deleted{deux \unsure{raffinements|raffinement} de}} \deleted{\unsure{$F$}}raffinement de $F$ l'un, de $K$ l'autre, et tels que $F'_K\overset{\unsure{\text{défn}|\text{dfn}}}{=}F'|K$ soit compatible avec $H'_K$. Si on peut en conclure qu'ils sont compatibles, on pourrait construire leur \unsure{Sup} $H'_F$ et $H'_F\longrightarrow F$, et symétriquement $H'_G$ (\unsure{Sup} de $G'$ et de $H'_K$). On vérifierait sans mal $H'_F|K=H'_K$, et de $=$ $H'_G|K=H'_K$, donc par l'axiome de recollement, on \unsure{saurait} $H'_F\mathrel{>\mkern-14mu<}H'_G$ et on pourrait construire $H'=H'_F\vee H'_G$. \deleted{\unsure{D}}En\unsure{.} résumé, pour vérifier la partie non triviale de la proposition, on est ramené au cas particulier où
\diagram{\begin{tikzcd}[ampersand replacement=\&,row sep=1.25em,column sep=0.25em]
G \arrow[d,phantom,"{\rotatebox{90}{$\ll$}}" description] \& \deleted{\unsure{\leq}}\unsure{\subset|<} \& F \arrow[d,phantom,"{\rotatebox{90}{$\ll$}}" description] \& \text{et où}\quad F'_G\leq G' \\
G' \& \& F' \&
\end{tikzcd}}
\scanpage{77}\mspage{7\deleted{\unsure{5}}6}
i.e. on est ramené au

\uline{Corollaire} Soient $G\leq F$, $F'\ll F$ un raffinement de $F$, induisant le raffinement $F'_G\deleted{\unsure{\ll}}\leq F'$ de $G$, et \deleted{soit} supposons $F'_G\leq G'$, où $G'$ est un raffinement de $G$. alors \deleted{$G'$} $F'$ et $G'$ sont compatibles (et $F'\cap G'=F'_G$) \deleted{$\unsure{F'\cap G'|F'\vee G'}\unsure{=|-}$}

Est-ce que ça résulte des autres axiomes ? \deleted{\unsure{Regardons}} On peut encore réduire plus (c'est \unsure{standard|standart}) de façon \unsure{à|:} obtenir seulement

\deleted{\uline{\unsure{At 12}}} \inserted{Corollaire} Soient
\diagram{\begin{tikzcd}[ampersand replacement=\&,column sep=0.4em,row sep=1.3em]
X \arrow[d,phantom,"{\rotatebox{90}{$\ll^{\unsure{\circ|\bullet}}$}}" description] \& > \& Y \arrow[d,phantom,"{\rotatebox{90}{$\ll$}}" description] \arrow[dr,phantom,"{\rotatebox{-45}{$\gg^{\circ}$}}" description] \& \\
X' \& \geq \& X'_{\unsure{Y|y}} \arrow[d,phantom,"{\rotatebox{90}{$=$}}" description] \& Y'\\
\& \& {\unsure{x|X}}'|Y \&
\end{tikzcd}}
Si $X'_Y\deleted{\unsure{\leq|<}}\mathrel{>\mkern-14mu<}Y'$, alors $X'\mathrel{>\mkern-14mu<}Y'$

C'est \inserted{(compte tenu de At 11 ou At $11'$)} un renforcement considérable de At 10, mais qui résulterait aussi de At $9'$ $+$ At 10, ou de At $9''$ $+$ At 10.

On va formuler plutôt, sous forme un peu plus générale qui va impliquer At $11'$ (donc, \unsure{en retour\,!|en retour}, aussi At 11) :

\uline{At 12} Soient $X'\ll X$ et $Y'\ll Y$ dans $M$, \deleted{tels que} avec $X\mathrel{>\mkern-14mu<}Y\unsure{\mathord{.}|\circ}$ Soit $K=X\cap Y$, d'où $X'_K\leq X'$, $X'_K\ll K$ et $Y'_K\leq Y'$, $Y'_K\ll K$. \unsure{Ceci|Ceci,}
\scanpage{78}\mspage{7\deleted{\unsure{6}}7}
posé, si $X'_K\mathrel{>\mkern-14mu<}Y'_K$, alors $X'\mathrel{>\mkern-14mu<}Y'$. En d'autres termes, si $\forall\, X''\in\widetilde{X'}$ \unsure{tels|tel} que $X''\ll K$ ($\Leftrightarrow X''\ll Y$), et tout $Y''\in\widetilde{Y'}$ \unsure{tels|tel} que $Y''\ll K$ ($\Longleftrightarrow Y''\ll X$), on a $X''\mathrel{>\mkern-14mu<}Y''$, alors $X'\mathrel{>\mkern-14mu<}Y'$.

Ceci implique At $11'$, \inserted{\deleted{\unsure{At 11}}} At 10, et la prop. de la page 73.
% slanted
\margin{\unsure{donc on peut} faire \unsure{sauter} At 10,}
Donc cela implique At 11 pour \deleted{familles \uline{finies}} $\widetilde F$ \uline{fini}, et le cor. 2 page 69 \unsure{pour} \unsure{ens|ens.} d'indices \uline{fini}.

\deleted{\unsure{Nous avons} un énoncé qui implique At 11} Les axiomes At 9 (critère de disjonction par les lieux) \inserted{\unsure{ses variantes} $9',9''$ (critères de compatibilité \inserted{par les lieux})} At 10 \unsure{.|'} (\inserted{premier} critère de disjonction des lieux), At 11 (critère de recollement des raffinements), \inserted{sa variante un peu affaiblie At $11'$,} \unsure{A12|At 12} (critère de compatibilité par intersection) sont de la nature de critères d'existence de figures (figures sommes, \deleted{\unsure{et}} en les quatre cas) intimement reliés entre eux, par des implications diverses\unsure{,|;} qu'on va résumer \unsure{des|dans} diagrammes. Mais je vais \inserted{d'abord} donner une formulation commune \deleted{\unsure{à}} \deleted{\unsure{qui|que}} à At $9'$ et At $9''$, par la \unsure{notion|création} de $\Sigma$-Atelier, où $\Sigma$ est une sous-cat. pleine de la cat. des ens. ordonnés. Un atelier est un $\Sigma$-atelier si
\[
\text{(i)}\quad \forall\,\deleted{\unsure{F}}\unsure{X|F}\in\deleted{\unsure{\mathcal F|\mathscr F}}M\qquad \widetilde{\deleted{F}X}\in\Sigma
\]
(ii) $\mathcal F$ satisfait l'axiome qui suit :
\scanpage{79}\mspage{7\deleted{\unsure{7}}8}
At $9_{\Sigma}$ (critère de compatibilité par les lieux, pour $\Sigma$-ateliers) Soient $X,Y\in M$, \deleted{tels que} $\deleted{\unsure{K|X}}K=X\cap Y$ ($=\operatorname{Inf}^{\leq}(\unsure{x|X},\unsure{Y|y})=\operatorname{Sup}_{Z\in\widetilde X\cap\widetilde Y} Z$), de sorte que $\widetilde K=\widetilde X\cap\widetilde Y$. Pour que $x\deleted{\unsure{\text{et}|\parallel}}\mathrel{>\mkern-14mu<}y$, \deleted{soient disjoints} \inserted{\unsure{on|en}} il \deleted{\unsure{(}}f. et \deleted{\unsure{)}}s. que
\[
\begin{aligned}
\text{(i)}\quad &\forall\,x\in\operatorname{omb}(X)\setminus\operatorname{omb}(K),\quad y\in\unsure{\operatorname{om}|\operatorname{omb}}(Y)\setminus\operatorname{omb}(K),\\
&\text{on a}\quad x\parallel y\quad\text{(ou encore }x\mathrel{>\mkern-14mu<}y\text{)}\\
\text{(ii)}\quad &\widetilde K\in\Sigma
\end{aligned}
\]
Cela implique donc

\uline{Corollaire 1} (ancien At $9''$) Si on a (i) et (ii$'$) $K\in M$ (\unsure{i.e|i.e.} $\widetilde K$ a un plus grand élément), alors $X\mathrel{>\mkern-14mu<}Y$.

En effet, \unsure{.} $K\in M\Longrightarrow\widetilde K\in\Sigma$.

\uline{Corollaire 2} (ancien At $9'$) Supposons que $\Sigma\supset$ \deleted{\unsure{cat.}}catégorie des \unsure{ens|ens.} ordonnés isomorphes à un $\widetilde F$, $F\in\unsure{\mathcal F|\mathscr F}$ (p.ex. $\Sigma={}$cat. \unsure{de|des} tous les ens. ordonnés) alors $X\mathrel{>\mkern-14mu<}Y$ sss on a (i).

% slanted
\margin{\unsure{et|en} \unsure{prenant} l'un ou l'autre cas : 1$^{\mathrm o}$) $\Sigma={}$ tous les \unsure{ens.|ens} ordonnés, \unsure{donc} (ii) est vide, 2$^{\mathrm o}$) $\Sigma={}$ ss-cat. pleine des ens \unsure{avec} plus grand élément. (de sorte que (ii) équivaut à $K\in M$.)}

\uline{Scholie} \deleted{On supposera que tout atelier} \inserted{Tous les ateliers auxquels j'ai songé} \deleted{est un} \inserted{sont des} $\Sigma$-ateliers pour $\Sigma$ convenable, \deleted{ce qui \deleted{\unsure{implique que}} \inserted{\deleted{\unsure{équivaut}}} à dire que c'est un $\Sigma$-atelier pour $\Sigma\unsure{.}=\Sigma_M={}$cat. des ens. ord. équivalents \unsure{à|:} un $\widetilde X$, $X\in M$. Cela signifie aussi que At $9_{\Sigma}$ est valable pour $\Sigma_M$, ou encore que l'on veut \uline{\unsure{préciser les univers}} de validité du cor. 1}
\scanpage{80}\mspage{7\deleted{\unsure{8}}9}
Ceci posé, on a

\diagram{\resizebox{\linewidth}{!}{\begin{tikzcd}[ampersand replacement=\&,column sep=1em,row sep=1.5em]
{} \& \mathrm{At}\,9_{\Sigma}\arrow[r,Rightarrow]\arrow[d,Rightarrow] \& \mathrm{At}\,9
\& \mathrm{At}\,11\arrow[r,Leftrightarrow]\arrow[d,Rightarrow,start anchor=south west] \& \begin{gathered}\mathrm{At}\,11\ \text{pour}\\ \deleted{\unsure{\mathcal F|\mathscr F|F}}F=\operatorname{Sup} F_i\\ \text{\unsure{quelc.|quelq.}}\end{gathered}\arrow[d,Rightarrow] \\
\deleted{\unsure{A|\mathrm{At}}} \& \mathrm{At}\,12\arrow[r,Rightarrow]\arrow[dr,Rightarrow,rounded corners=8pt,to path={|- (\tikztotarget) \tikztonodes}] \& \mathrm{At}\,10
\& \mathrm{At}\,11'\arrow[r,Leftrightarrow]\arrow[dl,Rightarrow] \& \begin{gathered}\mathrm{At}\,11\ \text{pour }F=\operatorname{Sup} F_{\unsure{i|I}}\\ \text{fini}\end{gathered} \\
{} \& {} \& \mathrm{At}\,11\deleted{\unsure{'}}\ \text{pour }\widetilde F\ \text{fini}
\& \deleted{\Rightarrow}\ \deleted{\unsure{\mathrm{At}\,11'|\mathrm{At}\,11}} \& {}
\end{tikzcd}}}

Dans ces axiomes, le seul qui \deleted{\unsure{f}} soit de\unsure{,} la nature d'un \deleted{\unsure{énoncé|axiome}} \inserted{critère} d'existence pour un \deleted{\unsure{sommet|somme}} Sup d'une famille \uline{infinie} de figures \deleted{\unsure{(ou d'un|(ou d'une}} est At 11 (qui pour cette raison n'est pas coiffé par At 12) -- le seul autre qui fasse \unsure{intervenir} un sup pour un $\unsure{\exists|J}$

\diagram{\resizebox{\linewidth}{!}{\begin{tikzcd}[ampersand replacement=\&,column sep=1em,row sep=1.4em,execute at end picture={
\draw[thin] ([xshift=1ex,yshift=1ex]a.north west) -- ([xshift=2ex,yshift=-1ex]d.south east);
\draw[thin] ([xshift=1.2em,yshift=2ex]a.north east) -- ([xshift=.5em,yshift=-2ex]c.south);
\draw[thin] (b.north west) -- (h.south west);
\draw[thin] (e.north) -- (d.south west);
\draw[thin] (e.south) -- (d.north east) -- ([xshift=.4em]e.south west) -- (d.north);
\draw[thin] (e.north east) -- (h.east);
\draw[thin] (g.north) -- (h.south west);
\draw[thin] (j.north east) -- (h.south east);
}]
|[alias=a]|\deleted{\mathrm{At}\,9_{\Sigma}}\arrow[r,Rightarrow]\arrow[d,Rightarrow] \& |[alias=b]|\deleted{\mathrm{At}\,9} \& |[alias=e]|\deleted{\mathrm{At}\,9_{\Sigma}+\mathrm{At}\,10}\arrow[dll,Rightarrow]\arrow[ddl,Rightarrow] \& {} \& {} \\
|[alias=c]|\deleted{\mathrm{At}\,12}\arrow[r,Rightarrow]\arrow[dr,Rightarrow,rounded corners=7pt,to path={|- (\tikztotarget) \tikztonodes}] \& \deleted{\mathrm{At}\,10} \& \deleted{\mathrm{At}\,11}\arrow[dl,Rightarrow] \& |[alias=g]|\deleted{\left(\Longleftrightarrow\begin{gathered}\text{\unsure{Recoll.|Rgcll.} des}\\ \text{raffinements}\\ \text{des }F=\operatorname{Sup} F_i\end{gathered}\right)} \& {} \\
{} \& |[alias=d]|\deleted{\mathrm{At}\,11'} \& |[alias=h]|\deleted{\left(\Longleftrightarrow\begin{gathered}\text{Recollement}\\ \text{des raffinements}\\ \text{des }F=\operatorname{Sup}_{i\in I} F_i\end{gathered}\right)}
\& |[alias=j]|\deleted{\Rightarrow} \& {} \\
{} \& {} \& \deleted{\text{pour }I\text{ fini}} \& {} \& {}
\end{tikzcd}}}

\diagram{\resizebox{\linewidth}{!}{\begin{tikzcd}[ampersand replacement=\&,column sep=.7em,row sep=1.4em]
{} \& {} \& \mathrm{At}\,9_{\Sigma}+\mathrm{At}\,10\arrow[dl,Rightarrow]\arrow[d,Rightarrow] \& {} \& {} \\
{} \& \mathrm{At}\,9_{\Sigma}\arrow[dl,Rightarrow] \& \mathrm{At}\,12\arrow[dl,Rightarrow]\arrow[dr,Rightarrow] \& {} \& \mathrm{At}\,11\unsure{.}\ \begin{gathered}\left(\Longleftrightarrow\begin{gathered}\text{\unsure{Recollemt|Recollement} des}\\ \text{raff. des }F=\operatorname{Sup} F_i\end{gathered}\right)\\ (\Longleftrightarrow\mathrm{At}\,11\,\text{bis}+\mathrm{At}\,11')\end{gathered}\arrow[dl,Rightarrow] \\
\begin{gathered}\mathrm{At}\,9\\ \text{(critère de disj.}\\ \text{par les lieux)}\end{gathered} \& \mathrm{At}\,10 \& {} \& \mathrm{At}\,11'\ \left(\Longleftrightarrow\begin{gathered}\text{\unsure{Recollemt|Recollement} des raffinements}\\ \text{des }F=\operatorname{Sup} F_i\end{gathered}\right) \& {} \\
{} \& {} \& {} \& \text{pour }I\text{ fini} \& {}
\end{tikzcd}}}

Le seul axiome qui échappe à $\mathrm{At}\,9_{\Sigma}+\mathrm{At}\,10$ (équivalent à $\mathrm{At}\,9_{\Sigma}+\mathrm{At}\,12$) est At 11, qui implique un Sup infini. On pourrait alors \unsure{renforcer} At 11 par un énoncé \deleted{qui} qui soit équivalent, compte tenu de At $11'$ :
